\scanpage{40}
\mspage{31}
% archive sheet number: 40
% inv: text | top | On peut dès lors
On peut dès lors se demander quels axiomes à mettre sur $(M,\leq,\overset{\circ}{\ll})$, pour que la relation
% inv: display | upper third | star disjointness definition
\[
X\mathrel{|\circ|}Y\Longleftrightarrow\{X,Y\}\text{ non minoré \unsure{ds|dans} }M_{\overset{\circ}{\ll}}\tag*{(*)}
\]
% inv: text | upper middle | donne avec
 donne, avec $\leq$, $\overset{\circ}{\ll}$, un magasin.

% inv: text | upper middle | Les axiomes
Les axiomes Mag 1, Mag 2 (p. 22\unsure{?}, à \unsure{un .|un \textquotesingle} \unsure{près} avec les $'$) sont clairs, \deleted{Mag $3'$ \unsure{est}} sur $\ll$, $\overset{\circ}{\ll}$, \deleted{\unsure{\&}} Mag $3'$ est automatique. Mag $4'$ \deleted{\unsure{donne une condition supplémentaire}} \inserted{résulte de Mag $2'$ a !}

% inv: deletion | middle | Autre
\deleted{\unsure{Autre}}

% inv: text | middle | Proposition
\uline{Proposition}\quad Les magasins $(M,\leq,\overset{\circ}{\ll})$ (avec $\mathrel{|\circ|}$ définissable par \unsure{(*)} ci-dessus) sont caractérisés par les conditions ;

% inv: text | middle | Mag 1 prime
\uline{Mag $1'$}\quad $\leq$, $\overset{\circ}{\ll}$ des relations d'ordre

% inv: text | middle | Mag 2 prime
\uline{Mag $2'$}\quad Comme l'homologue p. 22.

% inv: deletion | middle lower | Mag 3 prime
\deleted{\uline{Mag $3'$}}\quad\deleted{$\unsure{\forall X,Y,\ X\overset{\circ}{\ll}Y}$}

% inv: text | lower middle | C'est pourtant aussi simple
C'est \unsure{pourtant} aussi simple que le cas ensembliste\inserted{(p. 9)} ! Il n'y a pas à se mettre en frais d'un \unsure{Magquens}, comme page 24. Finalement, on a une hiérarchie de structures de plus en plus particulières, des \unsure{noti-magasins|anti-magasins} :

% inv: diagram | lower third | hierarchy of magasins
\diagram{\scriptsize\begin{tikzcd}[ampersand replacement=\&,column sep=small,row sep=small]
\& {}\arrow[dd,phantom,"{\left\{\vphantom{\begin{matrix}a\\b\\c\end{matrix}}\right.}"{anchor=center}] \& \text{Magasins généraux} \& \leq,\overset{\circ}{\ll},\mathrel{|\circ|}\quad\substack{\text{Mag 1--Mag 4}\\\text{(p. 22)}}\\
\substack{M\text{ est l'ens}\\\text{des bornes}\\\text{\unsure{principales|princip.}}} \& \& \text{Magasins }\leq,\overset{\circ}{\ll}\arrow[u,Rightarrow] \& \substack{\text{Mag }1',\text{Mag }2'=\text{Mag 2}\\\text{(cette page)}}\\
\& {} \& \text{Magasins quasi-ensemblistes}\arrow[u,Rightarrow] \& \text{Mag }1',\text{Mag }2',\text{Mag }3'\quad\text{(p. 24)}\\
\substack{L\text{ est l'ens des}\\\text{bornes princip.}} \& \{ \& \text{Magasins ensemblistes}\arrow[u,Rightarrow] \& \unsure{\mathrm{Magens}}1,\mathrm{Magens}2\quad\text{(p. 9)}
\end{tikzcd}}

% inv: text | bottom | On sera amené aussi
On sera amené aussi à regarder les magasins $\leq$, $\mathrel{|\circ|}$, qui auront d'autres avantages (qui se trouvent entre les magasins généraux, et les ensemblistes\ldots).
